\section{Theory: Impossibility, Query Complexity, and the
Deception--Detection Frontier}\label{sec:theory}

Three results organize the theory: passive monitoring cannot tomograph, so
measurement must be active; active estimation has a priced sample complexity,
and adaptivity pays only where propositions screen out; and a
patch-deception-decay law and an onset dichotomy jointly cap what deception
can buy and hide. Proofs and derivation-level commentary are in the
supplement; the inequality machinery is standard, and the
contribution is the objects it bounds.

\subsection{Passive monitoring cannot tomograph}\label{sec:t1}

\begin{theorem}[Passive-observation impossibility]\label{thm:passive}
There exist belief hypotheses $B_0\neq B_1$ and an adversary-controlled event
$E$ (e.g., a forged-authority directive) with $\Pr_\mu(E)=0$ such that:
\begin{enumerate}[label=(\roman*),leftmargin=2em]
\item the passive transcript distributions coincide,
      $\TV\bigl(P^{\mathrm{pas}}_{B_0},P^{\mathrm{pas}}_{B_1}\bigr)=0$; and
\item a single probe family separates them completely,
      $\E_{B_1}[X_B(\Pi)]-\E_{B_0}[X_B(\Pi)]=1$.
\end{enumerate}
\end{theorem}

\begin{corollary}[Detection impossibility under bounded false
positives]\label{cor:detection}
Any monitor whose input is the passive transcript stream---any classifier, any
alarm rule---has identical output distributions under $B_0$ and $B_1$. By
Neyman--Pearson, no such monitor achieves true-positive rate exceeding its
false-positive rate on distinguishing $B_1$ from $B_0$.
\end{corollary}

Two worlds share a transcript law and differ in belief: no passive statistic
separates them, and the active design does.

\subsection{Query complexity of active estimation}\label{sec:qc}

\begin{theorem}[Uniform estimation guarantee]\label{thm:uniform}
Fix accuracy $\varepsilon>0$ and confidence level $\alpha\in(0,1)$. If each
proposition $i$ receives
\[
n_i \;\ge\; \frac{1}{2\varepsilon^{2}}\,\ln\frac{2K}{\alpha}
\]
probes, then the empirical estimator $\hat\beta_i$ satisfies
\[
\Pr\Bigl[\,\max_{i\in[K]}\,\bigl|\hat\beta_i-\beta_i\bigr|\le\varepsilon
\Bigr]\;\ge\;1-\alpha
\]
(Hoeffding's inequality \cite{hoeffding1963} plus a union bound over the $K$
propositions).
\end{theorem}

\begin{proposition}[Adaptive screening gain for the certificate
objective]\label{prop:screening}
If at most $s\ll K$ propositions have $\beta_i\ge\tau$ for a threshold
$\tau>\varepsilon$, a two-phase policy---uniform screening at a coarse
width to identify above-threshold candidates, followed by
$\varepsilon$-refinement on the surviving set---achieves the certificate
decision of Section~\ref{sec:cert} with total probes
\[
O\bigl(K\,\tau^{-2}\ln(K/\alpha)+s\,\varepsilon^{-2}\ln(K/\alpha)\bigr),
\]
a saving concentrated on the $K-s$ below-threshold propositions whenever
$\varepsilon\ll\tau$. The uniform $\varepsilon$-guarantee of
Theorem~\ref{thm:uniform} itself admits no adaptive saving: it requires
$\varepsilon$-width on every proposition.
\end{proposition}

The uniform guarantee is tight in $\varepsilon$; adaptive screening pays only
when the deployment lets the design screen families out.

\subsection{The deception--detection frontier}\label{sec:t5}

\begin{definition}[Sensitivity]\label{def:sens}
The functional $\beta$ is \emph{$L$-sensitive} if for all transcript
distributions $\mu_1,\mu_2$ and propositions $\pi$,
\[
\bigl|\beta[\mu_1](\pi)-\beta[\mu_2](\pi)\bigr|
\;\le\;L\,\TV(\mu_1,\mu_2);
\]
the frontier caps below apply with the sensitivity constant so defined.
\end{definition}

\begin{remark}[Sensitivity is an identity for binary verdicts, not a
regularity condition]\label{rem:sens}\ \
For verdict channels $X_B\in\{0,1\}$ (Definition~\ref{def:beta}), sensitivity
is \emph{not} an empirical regularity condition that a protocol could
falsify: with $A_\pi=\{X_B(\Pi)=1\}$ the measured belief is
$\beta[\mu](\pi)=\mu(A_\pi)$, so
\[
\begin{aligned}
\bigl|\beta[\mu_1](\pi)-\beta[\mu_2](\pi)\bigr|
&=|\mu_1(A_\pi)-\mu_2(A_\pi)|\\
&\le\sup_B|\mu_1(B)-\mu_2(B)|\\
&=\TV(\mu_1,\mu_2).
\end{aligned}
\]
i.e.\ \emph{$L=1$ always}, by the sup-over-events definition of $\TV$.
The dose--response experiment E26 (Section~\ref{sec:limits}) instantiates
this identity on the constructed planes---by construction, since the
injected dose shift is a $\TV$ shift and the verdict is binary---and
therefore cannot return a falsifying value; its confirmation counts as a
consistency check of the harness, not a test of an assumption. The frontier
caps in this section consequently carry $L=1$ and are tight in that
constant; sensitivity over \emph{probe} space (a Lipschitz condition on the
proposition $\pi$ rather than the transcript distribution) is a strictly
stronger property that does \emph{not} follow from the identity, is
violated by the probe-aware constructions of
Section~\ref{sec:surface}, and is assumed explicitly---as an
assumption---only where the transfer result of
Proposition~\ref{prop:gaming} needs it.
\end{remark}

\begin{lemma}[Amplification]\label{lem:amp}
Let $\delta=\TV(P,Q)\in(0,1]$ and let $P^N,Q^N$ denote the distributions of
$N$ i.i.d.\ samples from $P$ and $Q$. Then
\[
\TV\bigl(P^{N},Q^{N}\bigr)\;\ge\;1-2\exp\!\bigl(-N\delta^{2}/2\bigr).
\]
\end{lemma}

\begin{theorem}[Patch-deception decay]\label{thm:patch}
Let $\beta$ be $L$-sensitive and let a deception perturb attacker-accessible
observations with $\delta=\TV(P,Q)$. Then:
\begin{enumerate}[label=(\roman*),leftmargin=2em]
\item every proposition's measured shift is footprint-bounded,
      $\;|\beta[Q](\pi)-\beta[P](\pi)|\le L\delta$ for every $\pi$;
\item an attacker with horizon $N$ who knows the honest distribution $P$ and
      uses the threshold test $\psi^{c}=\mathbf{1}[S<t]$ detects the
      deception with probability at least $1-2e^{-N\delta^{2}/2}$, and no
      test's power-minus-size exceeds $\TV(P^{N},Q^{N})$ (the attacker who
      must learn the honest rate from a memorized batch is the subject of
      Theorem~\ref{thm:onset});
\item holding detection probability at most $\gamma$ at horizon $N$ forces
      \[
      \begin{aligned}
      &\delta\le\sqrt{2\ln\bigl(2/(1-\gamma)\bigr)/N},\\
      &\text{hence gain}\ \le L\sqrt{2\ln\bigl(2/(1-\gamma)\bigr)/N}
      =O\bigl(L/\sqrt N\bigr).
      \end{aligned}
      \]
\end{enumerate}
\end{theorem}

\begin{theorem}[Onset dichotomy]\label{thm:onset}
Let the attacker sample the environment for $M$ turns and let deception
change the observable event rate from $p$ to $q=p+\delta$ (the registered
planes shift \emph{up}; the footprint remains $\delta>0$). Write
$\sigma_0=\sqrt{p(1-p)/N_0}$, $z=z_{1-\beta_d/2}$ and $z'=z_{1-\beta_d}$,
and define the two boundaries
\[
\rho_0:=z\,\sigma_0
\quad\text{and}\quad
\Delta_{\beta_d}:=(z+z')\sigma_0=1.84\,\rho_0
\]
at $\beta_d=0.05$. Here $\rho_0$ is the \emph{reference width} (a two-sided CI half-width of the
memorized batch---the level at which the reference \emph{resolves} the
deception) and the \emph{certified-detection boundary} $\Delta_{\beta_d}$
(the shift at which the level-$\beta_d$ two-sample $z$-test actually
\emph{attains} power $1-\beta_d$---at $\delta=\rho_0$ its power is
$\approx0.50$, not $1-\beta_d$). If deception activated mid-session and
the attacker memorized $N_0\ge1$ honest samples, three regimes hold:
\begin{enumerate}[label=(\roman*),leftmargin=2em]
\item \emph{Known onset.} One of three regimes holds:
  \begin{enumerate}[label=(\alph*),leftmargin=2em]
  \item \emph{Certified detection}, $\delta>\Delta_{\beta_d}$: the
        attacker detects with probability $\ge1-\beta_d$ once
        \[
        M\ge M^{*}(\delta):=
        \frac{(z+z')^{2}q(1-q)}
        {\delta^{2}-(z+z')^{2}p(1-p)/N_0},
        \]
        the smallest horizon at which certified power is achievable:
        the denominator is $\delta^{2}-\Delta_{\beta_d}^{2}$, so
        $M^{*}$ is finite iff $\delta>\Delta_{\beta_d}$.
        At the registered planes $(p,q)=(0.3,0.7)$, where both variances equal $0.21$,
        $M^{*}\approx N_0/\bigl((\delta/\Delta_{\beta_d})^{2}-1\bigr)$:
        $M^{*}=30$ at the $(N_0,\delta)=(40,0.4)$ cell---a certified
        horizon scaling as $\delta^{-2}$: asymptotically
        $M^{*}\approx(z+z')^{2}\bar v/\delta^{2}=2.73/\delta^{2}$, and
        at this cell the finite-reference factor
        $1/(1-(z+z')^{2}\bar v/(N_0\delta^{2}))=1.74$ raises the constant
        to $4.8$, giving $M^{*}\approx4.8/\delta^{2}\approx30$. No
        horizon certifies below the boundary: at $\delta=\rho_0$ the
        power is $\approx0.50$, and at the two middle registered cells
        ($\delta/\sigma_0=2.76<3.6$) the limiting power is
        $\Phi(0.80)=0.79$; discreteness does not close the gap, since the
        exact one-sided binomial test at level $\beta_d$ with $N_0$
        reference draws has power $0.850$ at $(10,0.4)$ and $0.785$ at
        $(40,0.2)$ (one-sided normal analogue $0.868$), and a test on
        those draws alone cannot grow with $M$.
  \item \emph{Uncertified but above the reference width},
        $\rho_0\le\delta\le\Delta_{\beta_d}$: the same test's power
        converges with $M$ to
        $\Phi\!\bigl(\delta/\sigma_0-z\bigr)<1-\beta_d$: \emph{no} horizon
        attains certified power, and deception remains visible (power
        bounded away from level) but never refutable at the registered
        level.
  \item \emph{Below the reference width}, $\delta<\rho_0$: no
        level-$\beta_d$ two-sample test of the fresh stream against the
        memorized reference whose rejection threshold is fixed
        independently of the estimated honest rate reaches power
        $1-\beta_d$ for any $M$: the threshold must clear the $p$-robust
        reference width. The gain is bounded by
        \[
        L\,\rho_0\;=\;L\,z\,\sqrt{p(1-p)/N_0}
        \;=\;O\bigl(L/\sqrt{N_0}\bigr).
        \]
  \end{enumerate}
  In all three regimes the deception's price is set by the attacker's
  memorized reference size $N_0$, not by dwell time: certified refutation
  at a finite horizon (a), saturated sub-certified detection (b), or a
  gain capped at the reference width (c).
\end{enumerate}
\end{theorem}

\begin{remark}[Birth onset is a channel identity]\label{rem:birth}
If the deception plane is active from the attacker's first packet
($N_0=0$), the attacker holds one sample path from a single unknown
distribution and has no in-channel reference. Call an \emph{in-channel}
test any measurable function of the transcript path alone, and a
deception's \emph{plausible set} the set of honest worlds it cannot be
distinguished from by in-channel evidence. Then ``deceptions inside the
plausible set are in-channel undetectable for every $M$'' is true
\emph{by the definition of the plausible set}---a channel identity, not a
falsifiable prediction, and we state it as one. Its content is
architectural: birth-onset endurance is governed by onset invisibility
and by the attacker's \emph{out-of-band} channels, not in-channel
sample size. A plane in an attacker's plausible set is not decidable
from within the channel; the
priors-over-worlds extension is open (Section~\ref{sec:discussion}); and
the E5 birth-onset readings of Section~\ref{sec:e5} are reported as
descriptive statistics, not as a test of this remark.
\end{remark}

\begin{corollary}[The deception--detection frontier]\label{cor:frontier}
For the detection statistics analyzed in Theorems~\ref{thm:patch}
and~\ref{thm:onset}---the threshold test on the amplifying event, and
level-$\beta_d$ two-sample tests against a memorized reference---achievable
deception gain is jointly capped at $O(L/\sqrt{N})$ and $O(L/\sqrt{N_0})$.
The object the analysis identifies is the frontier curve of achievable gain
against attacker probe budget $M$, out-of-band audit rate, and reference
size $N_0$; the adaptive minimax extension remains open
(Section~\ref{sec:discussion}).
\end{corollary}

The frontier composes the two caps: a patch buys at most $L\delta$ of shift,
stealth forces $\delta=O(\sqrt{\log(1/(1-\gamma))/N})$, and the onset dichotomy
locates a deception relative to the reference width.

\subsection{Empirical test: the onset structure (E5)}\label{sec:e5}
The registered E5 prediction---advantage $\to1$ above the reference floor and
flat at chance below it, with no grid---failed at three of four cells, and
the corrected theorem of Section~\ref{sec:t5} explains why: the two middle
cells ($\delta/\sigma_0=2.76<3.6$) cannot reach certified power at any
horizon, and the birth-onset self-test reads $0.431$--$0.529$ at
$\delta=0.2$ (two of four readings below chance) against a registered
``flat at chance''. E5 is exploratory in the program's split
(Section~\ref{sec:limits}). The frozen cells, the post hoc ordinal reading,
and the corrected-boundary table are in the supplement's E5 record.
